% ===== PRÉAMBULE CANONIQUE — à coller VERBATIM en tête de CHAQUE .tex =====
\documentclass[11pt,a4paper]{article}
\usepackage[utf8]{inputenc}\usepackage[T1]{fontenc}\usepackage{lmodern}
\usepackage[french]{babel}
\usepackage{amsmath,amssymb,mathtools}\usepackage{icomma}
\usepackage[version=4]{mhchem}          % \ce{...} : équations & formules
\usepackage{siunitx}\sisetup{locale=FR} % \qty{}{}, \si{}, \ang{}
\usepackage{chemfig}                    % \chemfig{...} : structures développées
\usepackage{xcolor}\usepackage{colortbl}
\usepackage{tikz}\usepackage{pgfplots}\pgfplotsset{compat=1.18}
\usepackage[most]{tcolorbox}\usepackage{enumitem}\usepackage{array,booktabs}
\usepackage[a4paper,margin=2.2cm]{geometry}
\usetikzlibrary{arrows.meta,calc,angles,quotes,positioning,patterns,backgrounds,shapes.geometric}
\definecolor{primary}{RGB}{222,49,99}\definecolor{primarydark}{RGB}{150,22,55}
\definecolor{defcol}{RGB}{0,114,178}\definecolor{methcol}{RGB}{52,92,148}
\definecolor{excol}{RGB}{0,135,90}\definecolor{attncol}{RGB}{200,40,55}
\tcbset{boxbase/.style={enhanced,breakable,boxrule=0.9pt,arc=2.5pt,
  left=3mm,right=3mm,top=2.3mm,bottom=2.3mm,fonttitle=\bfseries,coltitle=white}}
% --- boîtes TUTORIEL (noms EXACTS, ne pas renommer) ---
\newtcolorbox{cle}[1][]{boxbase,colback=defcol!6,colframe=defcol,
  title={\textsc{À retenir}\ifx\relax#1\relax\relax\else\textmd{ : \itshape #1}\fi}}
\newtcolorbox{methode}[1][]{boxbase,colback=methcol!7,colframe=methcol,sharp corners,
  title={\textsc{Méthode}\ifx\relax#1\relax\relax\else\textmd{ --- \itshape #1}\fi}}
\newtcolorbox{exemplebox}[1][]{boxbase,colback=excol!5,colframe=excol,
  title={\textsc{Exemple}\ifx\relax#1\relax\relax\else\textmd{ : \itshape #1}\fi}}
\newtcolorbox{piege}[1][]{boxbase,colback=attncol!6,colframe=attncol,
  title={\textsc{Piège}\ifx\relax#1\relax\relax\else\textmd{ : \itshape #1}\fi}}
\newtcolorbox{remarque}[1][]{boxbase,colback=black!4,colframe=black!45,
  title={\textsc{Remarque}\ifx\relax#1\relax\relax\else\textmd{ : \itshape #1}\fi}}
% --- boîtes EXERCICE / CORRIGÉ (noms EXACTS, auto-comptées) ---
\newtcolorbox[auto counter]{exo}[1][]{boxbase,colback=primary!4,colframe=primary,
  title={\textsc{Exercice}~\thetcbcounter\ifx\relax#1\relax\relax\else\textmd{ --- \itshape #1}\fi}}
\newtcolorbox[auto counter]{corrige}[1][]{boxbase,colback=excol!4,colframe=excol,
  title={\textsc{Corrigé}~\thetcbcounter\ifx\relax#1\relax\relax\else\textmd{ --- \itshape #1}\fi}}
\newcommand{\R}{\mathbb{R}}\newcommand{\N}{\mathbb{N}}\newcommand{\Z}{\mathbb{Z}}\newcommand{\ds}{\displaystyle}
\begin{document}
\begin{exo}[Les cinq visages du carbonyle]
Chacun de ces groupements contient un carbonyle $\ce{C=O}$. Nomme la fonction correspondante.
\begin{enumerate}[label=\alph*)]
  \item $\ce{R-CO-H}$
  \item $\ce{R-CO-}\mathrm{R}'$
  \item $\ce{R-CO-OH}$
  \item $\ce{R-CO-O-}\mathrm{R}'$
  \item $\ce{R-CO-NH2}$
\end{enumerate}
\end{exo}

\begin{exo}[Isomères de fonction]
On considère le propanal \ce{CH3-CH2-CHO} et la propanone \ce{CH3-CO-CH3}.
\begin{enumerate}[label=\alph*)]
  \item Donne la formule brute de chacun.
  \item Quelle fonction porte chacun ?
  \item Comment appelle-t-on deux molécules de même formule brute mais de fonctions différentes ?
\end{enumerate}
\end{exo}

\begin{exo}[Amine ou amide ?]
Nomme la fonction azotée de chaque molécule ; pour les amines, précise la classe.
\begin{enumerate}[label=\alph*)]
  \item \ce{CH3-CH2-NH2}
  \item \ce{CH3-CO-NH2}
  \item \ce{(CH3)3N}
\end{enumerate}
\end{exo}

\begin{exo}[Recenser les fonctions d'une molécule bifonctionnelle]
Liste \emph{toutes} les fonctions présentes dans chaque molécule.
\begin{enumerate}[label=\alph*)]
  \item \ce{HO-CH2-CH2-CHO}
  \item \ce{H2N-CH2-COOH} (glycine)
\end{enumerate}
\end{exo}

\begin{exo}[Classer les amines]
Comme les alcools, les amines se classent selon le nombre de carbones portés par l'azote. Donne la classe (primaire, secondaire ou tertiaire).
\begin{enumerate}[label=\alph*)]
  \item \ce{CH3-NH2}
  \item \ce{CH3-NH-CH3}
  \item \ce{(CH3)3N}
\end{enumerate}
\end{exo}

\end{document}
